% =====================================================================
%  CSE499A — EHR Pre-Consultation Medical Documentation System
%  Standalone section: Ethical and Professional Responsibility
%  Authors: Rafiur Rahman Mashrafi (2221971042),
%           M.G. Rabbi Hossen (2222516042),
%           Israt Zaman Srity (2211084042)
%  Advisor: Dr. Mohammad Ashrafuzzaman Khan [AzK]
%  Department of Electrical and Computer Engineering,
%  North South University · Spring 2026
% =====================================================================

\documentclass[12pt,a4paper]{article}

\usepackage[utf8]{inputenc}
\usepackage[T1]{fontenc}
\usepackage{mathptmx}
\usepackage[a4paper,margin=1in]{geometry}
\usepackage{setspace}
\usepackage{enumitem}
\usepackage{array}
\usepackage{tabularx}
\usepackage{booktabs}
\usepackage[hidelinks]{hyperref}
\usepackage{titlesec}
\usepackage{xcolor}
\usepackage{quoting}

\onehalfspacing
\setlength{\parskip}{6pt}
\setlength{\parindent}{1.5em}

\titleformat{\section}{\normalfont\large\bfseries}{\thesection}{0.7em}{}
\titleformat{\subsection}{\normalfont\normalsize\bfseries}{\thesubsection}{0.7em}{}

\title{\bfseries Ethical and Professional Responsibility:\\
EHR-Based Pre-Consultation Medical Documentation System}
\author{Rafiur Rahman Mashrafi (ID 2221971042)\\
        M.~G.~Rabbi Hossen (ID 2222516042)\\
        Israt Zaman Srity (ID 2211084042)\\[6pt]
        \emph{Faculty Advisor:} Dr. Mohammad Ashrafuzzaman Khan [AzK]\\
        Department of Electrical and Computer Engineering\\
        North South University, Bangladesh}
\date{Senior Design Project (CSE499A) \quad Spring 2026}

\begin{document}
\maketitle

\section{Introduction}

This section examines the ethical and professional responsibilities
that arise when building the
\textit{EHR-Based Pre-Consultation Medical Documentation System}.
The system handles three things at once that make ethical analysis
non-optional: \textit{sensitive personal health data} of identifiable
patients; \textit{language-dependent recognition} where errors fall
unevenly on speakers of less-resourced dialects; and \textit{a
clinical decision context} in which a misrecognised symptom or
medication could lead to a wrong consultation. The discussion below
is structured around the seven General Ethical Principles of the
\textit{ACM Code of Ethics and Professional Conduct}~\cite{eth:acm2018}.
For each principle, the project's design decisions are shown
explicitly --- not as aspirations, but as concrete actions or
boundaries already enforced in the system.

The team treats the ACM Code as the binding professional standard
for the software-engineering work of this project. The Code makes
clear that it is not a series of abstract slogans but a tool that
``the computing professional should consider in
ethical decision-making''~\cite{eth:acm2018}. The seven principles
in Section~1 of the Code (the General Ethical Principles) are
treated below as numbered duties of the project team.

\section{Mapping to the ACM Code of Ethics, Section 1}

\subsection{ACM 1.1 --- Contribute to society and to human well-being,
acknowledging that all people are stakeholders in computing}

The Code asks software professionals to ``promote fundamental human
well-being''~\cite{eth:acm2018} and to consider the effects of their
work on \emph{all} affected parties --- not only the immediate
purchaser. In the present project the well-being beneficiaries are
multiple: the patient (who gains a system that listens in their own
dialect), the physician (whose documentation burden is reduced), and
the wider Bangladeshi public-health system, which currently lacks
structured digital records in most small clinics.

The team has taken three concrete actions toward this principle.
First, the system is designed to assist patients who are
under-served by existing English-only digital health tools, in
particular rural and dialect-speaking populations. Second, it is
explicitly assistive rather than prescriptive: the AI does not
issue a diagnosis, and the physician's clinical authority is
unchanged. Third, the project's empirical findings (the dataset, the
benchmark scripts, the conclusion that small specialised models
outperform 7B-parameter audio LLMs on dialect Bangla) are
documented and shared in the public project repository so that they
contribute to the wider Bangla-NLP research community rather than
remaining locked inside one institution.

\subsection{ACM 1.2 --- Avoid harm}

The Code states that ``harm'' includes ``unjustified physical or
mental injury, unjustified destruction or disclosure of information,
and unjustified damage to property, reputation, and the
environment''~\cite{eth:acm2018}. The most plausible harm pathway in
this system is the propagation of an ASR or NER error into the
clinical record, leading the physician toward a wrong line of
inquiry. Three design choices mitigate this:

\begin{enumerate}[leftmargin=*]
\item \textbf{Original transcript visibility.} The verbatim ASR
transcript is always displayed alongside the structured record on
the doctor's dashboard. The physician can compare the two at a
glance and detect mis-extractions before they affect care.
\item \textbf{No auto-confirmed prescriptions or diagnoses.} The
system never recommends a specific drug, dose, or diagnosis. It
surfaces conditions worth ruling out, which is decision support, not
decision making.
\item \textbf{Honest accuracy reporting.} Section~4 of the project
report documents Word Error Rate, Character Error Rate, and BLEU on
real dialect speech, including the failure cases. The physician
sees an honest estimate of system reliability, not a marketing
figure derived from a clean read-aloud benchmark.
\end{enumerate}

The Code further requires that ``the person doing the work has a
responsibility to assess the risks and either mitigate them or
inform affected parties.''~\cite{eth:acm2018} The accuracy
limitations are explicitly named in the report's Limitations
section so that affected parties (instructor, future
implementers, and eventually clinicians) are informed.

\subsection{ACM 1.3 --- Be honest and trustworthy}

ACM 1.3 calls for transparency about ``system capabilities,
limitations, and potential problems.''~\cite{eth:acm2018} The team
has applied this in three places.

\textit{In the model evaluation:} the report does not cherry-pick
favourable models. All twelve baseline ASR models are reported,
including the Whisper-Small case where Word Error Rate exceeds
195\% (the model hallucinates more text than it transcribes). The
two larger models that failed entirely (Phi-4 Multimodal,
Qwen2.5-Omni) are reported as failures rather than excluded.

\textit{In the description of project status:} the report
distinguishes clearly between phases that are complete (Phases 1--3,
the dataset and the two ASR benchmarks), a phase that is in progress
as a \emph{research} activity (Phase~4, fine-tuning methodology
study --- no training has actually been performed), and a phase that
is selected but not yet executed (Phase~5, the chatbot comparison).
The team rejects the temptation to imply more work has been done
than is the case.

\textit{In the description of fitness for clinical use:} the
report does not claim the system is currently safe for clinical
deployment. It is a Senior Design~I research artefact whose
strongest baseline still mis-transcribes nearly half the words on
real dialect audio. Any deployment claim would be premature, and
the team does not make one.

\subsection{ACM 1.4 --- Be fair and take action not to discriminate}

ACM 1.4 requires that ``technologies and practices should be as
inclusive and accessible as possible.''~\cite{eth:acm2018} For an
ASR system, fairness is concretely measurable: does the model
perform comparably across dialects, genders, age groups, and
recording conditions? An ASR system that performs well on the speech
of urban men in their twenties but poorly on the speech of rural
older women would amplify existing healthcare inequalities rather
than reduce them.

The project addresses this in three ways:

\begin{enumerate}[leftmargin=*]
\item \textbf{Multi-dialect data collection.} The Phase~1 corpus
covers five regional varieties (Puran Dhaka, Barishali, Sylheti,
Normal Bangla, Indian Bangla) rather than collapsing all input to a
single ``standard'' Bangla. The Limitations section names two
varieties (Chittagonian, Rangpuri) that are currently
under-represented, scheduled for inclusion in CSE499B.
\item \textbf{Demographic logging.} The dataset log records speaker
gender and approximate age for every audio file. This makes
per-group accuracy analysis possible and was an explicit decision
to enable later fairness audits.
\item \textbf{Voice-first by design.} The interface does not assume
literacy, English fluency, or smartphone fluency. The Code's
inclusion principle is operationalised at the interaction layer,
not just at the model layer.
\end{enumerate}

The team is aware that data balancing alone does not eliminate bias
and that a true fairness audit will require larger per-group sample
sizes than CSE499A's 4.7-hour corpus permits.

\subsection{ACM 1.5 --- Respect the work required to produce new
ideas, inventions, creative works, and computing artifacts}

ACM 1.5 covers proper attribution and respect for intellectual
property. The project follows this in two specific practices.

First, every model evaluated in the project is used under its
publicly stated license, and the model authors are cited in the
project's references. Pretrained checkpoints are downloaded from
Hugging~Face's official repositories for Whisper, MMS, Wav2Vec~2.0,
SeamlessM4T, the Qwen2/3-Audio family, BengaliAI, Vakyansh, Voxtral,
and Phi-4 Multimodal --- not redistributed or modified silently.

Second, the open-source software stack used in the project (PyTorch,
Hugging~Face Transformers, Librosa, FFmpeg, WhisperX, jiwer, yt-dlp)
is acknowledged in the report's Methodology chapter. The project's
own contributions --- the dataset log, the evaluation scripts, and
the comparative results CSV --- are released back into the public
repository so that future students and researchers can build on
this work without re-deriving it.

\subsection{ACM 1.6 --- Respect privacy}

ACM 1.6 states that ``computing professionals should only use
personal information for legitimate ends and without violating the
rights of individuals and groups.''~\cite{eth:acm2018} Audio
recordings of patient speech are among the most sensitive personal
data that exist: voice is a biometric identifier; health status is
sensitive information under any reasonable data-protection regime.

The project enforces six privacy controls, of which the first
three apply already to the CSE499A research dataset and the
remaining three are deployment-time requirements specified in the
system design:

\begin{enumerate}[leftmargin=*]
\item \textbf{No real patient data in CSE499A.} The 4.7-hour Phase~1
corpus consists of publicly available recordings, not real clinical
encounters. No identifiable Bangladeshi patient's audio has been
captured for this project.
\item \textbf{Pseudonymous file naming.} The dataset log uses a
fixed pattern of the form
\textit{[dialect-code]\_[index]\_[gender]\_[age].mp3}
rather than recording any speaker identity.
\item \textbf{Restricted repository visibility for raw audio.} The
GitHub project repository contains code, the dataset log, and
evaluation results, but the raw audio files are kept in a
controlled-access Google Drive folder rather than embedded in the
public repository.
\item \textbf{Informed consent at deployment.} The deployed system
will require explicit, recorded informed consent in the patient's
own language before any audio capture. Patients can opt out at any
time without penalty.
\item \textbf{Encryption at rest and in transit.} Audio blobs and
EHR records will be encrypted both at rest and over the network,
and access to the dashboard will be gated by role-based access
control.
\item \textbf{Data sovereignty.} Wherever feasible the deployed
system will keep sensitive Bangladeshi health data on
infrastructure located within Bangladesh rather than on foreign
cloud providers.
\end{enumerate}

\subsection{ACM 1.7 --- Honor confidentiality}

ACM 1.7 requires that confidentiality be honored ``except in cases
where it is evidence of the violation of law''~\cite{eth:acm2018}.
For this project the principle applies to two distinct types of
confidential information.

\textit{Patient-level confidentiality:} the EHR record produced by
the pipeline is accessible only to the licensed physician and to
authorised clinical staff under role-based access control. It is
never shared with third parties for marketing, training, or
analytics purposes without separate, explicit, opt-in consent.

\textit{Project-level confidentiality:} the source audio used in
CSE499A is from publicly available recordings, but the team treats
even this as data to be handled responsibly --- not redistributed
casually, not used for purposes other than the stated research, and
clearly attributed to its source in the dataset log.

\section{Mapping to the ACM Code of Ethics, Section 2:
Professional Responsibilities}

Section~2 of the ACM Code lists professional responsibilities that
extend the General Principles to working software professionals.
Three of these are particularly load-bearing for the project.

\subsection{ACM 2.5 --- Give comprehensive and thorough evaluations
of computer systems and their impacts, including analysis of
possible risks}

The project's own internal evaluation discipline is designed
precisely to satisfy this principle. The Phase~2 and Phase~3
benchmarks were comprehensive: twelve baseline ASR models and six
larger multimodal audio LLMs were each evaluated across three
metrics (WER, CER, BLEU) on the same multi-dialect corpus, and the
results --- including failures --- were documented. The risk
analysis in this section names the harm pathways (mis-extraction,
fairness failure, privacy breach) explicitly.

\subsection{ACM 2.6 --- Perform work only in areas of competence}

ACM 2.6 cautions professionals to ``perform work only in areas of
competence''~\cite{eth:acm2018}. The team is explicit that medical
diagnosis is \emph{not} the team's area of competence. This is why
the system never issues a diagnosis: that judgment is reserved for
the licensed physician. The team's areas of competence ---
machine-learning engineering, speech recognition evaluation,
software-system design --- are exactly the parts of the system that
the team builds and reports on. The differential-diagnosis
\emph{hint layer} that surfaces conditions worth ruling out will be
constructed in CSE499B in consultation with clinical advisors, not
authored by the engineering team alone.

\subsection{ACM 2.9 --- Design and implement systems that are
robustly and usably secure}

ACM 2.9 lists security as a professional duty. The system design
specifies authentication (JWT), role-based access control,
encryption at rest and in transit, audit logging on the dashboard,
and pseudonymous identifiers in the audio store. The deployment
plan additionally calls for keeping audio and EHR records on
infrastructure located within Bangladesh wherever feasible, which is
a security as well as a sovereignty consideration.

\section{Other Applicable Codes and Standards}

Beyond the ACM Code, the project also takes guidance from:

\begin{itemize}[leftmargin=*]
\item the \textbf{IEEE Code of
Ethics}~\cite{eth:ieee}, in particular its first commitment ``to
hold paramount the safety, health, and welfare of the public,''
which is consistent with the project's decision to keep the
physician as the final clinical authority;
\item the \textbf{Software Engineering Code of Ethics}
(IEEE-CS / ACM)~\cite{eth:secode}, with its eight principles
covering public interest, client/employer, product, judgement,
management, profession, colleagues, and self;
\item the \textbf{WHO ``Ethics and Governance of AI for
Health''} guidance~\cite{eth:who}, whose six principles --- protect
autonomy, promote human well-being, ensure transparency,
foster responsibility, ensure inclusiveness, promote sustainable AI
--- align directly with the design choices already enumerated in
this section.
\end{itemize}

The project also anticipates compliance with the emerging Bangladesh
data-protection framework when the system is moved from research
artefact to clinical deployment.

\section{Conclusion}

The ACM Code of Ethics is not a checklist that one applies once at
the end of a project; it is, as the Code itself
states~\cite{eth:acm2018}, a tool for ongoing ethical
decision-making throughout development. The team has tried to use
it that way. Each of the seven General Ethical Principles maps onto
specific design decisions the project has already made:
sharing of dataset and evaluation artefacts (1.1), transcript
visibility and no auto-confirmed prescriptions (1.2), honest
reporting of model failures and project status (1.3), multi-dialect
coverage with demographic logging (1.4), proper attribution of all
models and tools used (1.5), six concrete privacy controls (1.6),
and role-based confidentiality both at the patient level and at the
project level (1.7). The professional responsibilities of Section~2
of the Code ---  comprehensive evaluation, working within
competence, building robustly secure systems --- are similarly
operationalised in the empirical chapters of the project report and
in the planned deployment-layer work for CSE499B.

The most important ethical commitment, restated from the system
design and worth repeating in this section, is the assistive
boundary: the AI helps the physician; the AI does not replace the
physician. The licensed clinician remains the final and
accountable decision-maker at every step of the consultation. The
team believes that holding to this boundary --- and being honest
about the system's current limitations --- is the most
professionally responsible position consistent with the ACM Code.

% =====================================================================
% REFERENCES
% =====================================================================
\begin{thebibliography}{99}

\bibitem{eth:acm2018}
ACM Council,
``ACM Code of Ethics and Professional Conduct,'' 2018.
\url{https://www.acm.org/code-of-ethics}.

\bibitem{eth:ieee}
IEEE,
``IEEE Code of Ethics,'' 2020.
\url{https://www.ieee.org/about/corporate/governance/p7-8.html}.

\bibitem{eth:secode}
IEEE-CS / ACM Joint Task Force on Software Engineering Ethics,
``Software Engineering Code of Ethics and Professional Practice,''
Version 5.2, 1999.
\url{https://ethics.acm.org/code-of-ethics/software-engineering-code/}.

\bibitem{eth:who}
World Health Organization,
``Ethics and Governance of Artificial Intelligence for Health:
WHO Guidance,'' 2021.
\url{https://www.who.int/publications/i/item/9789240029200}.

\bibitem{eth:beauchamp}
T.~L.~Beauchamp and J.~F.~Childress,
\emph{Principles of Biomedical Ethics}, 8th ed.,
Oxford University Press, 2019.

\bibitem{eth:char}
D.~S.~Char, N.~H.~Shah, and D.~Magnus,
``Implementing Machine Learning in Health Care --- Addressing
Ethical Challenges,''
\emph{New England Journal of Medicine}, vol.~378, pp.~981--983,
2018.

\bibitem{eth:vayena}
E.~Vayena, A.~Blasimme, and I.~G.~Cohen,
``Machine Learning in Medicine: Addressing Ethical Challenges,''
\emph{PLOS Medicine}, vol.~15, no.~11, e1002689, 2018.

\bibitem{eth:gdpr}
European Parliament and Council of the European Union,
``Regulation (EU) 2016/679 (General Data Protection Regulation),''
2016.
\url{https://gdpr-info.eu}.

\end{thebibliography}

\end{document}
